% Half-title, title and copyright pages.
%
% \bookversion is the release the PDF was built for: `make VERSION=book-v2` passes it in, and the
% release workflow passes the tag. Without it the title page carries only the date of the build.
% Starred, so that an undefined version compares equal to \empty below, as an empty one does.
\providecommand*{\bookversion}{}
\newcommand{\bookdate}{\number\day\space\ifcase\month\or January\or February\or March\or April\or
  May\or June\or July\or August\or September\or October\or November\or December\fi\space
  \number\year}
\newcommand{\bookedition}{\ifx\bookversion\empty\else\bookversion\,·\,\fi\bookdate}

\thispagestyle{empty}
\vspace*{2.2in}
\begin{flushright}
  {\Large\bfseries Machine Learning}
\end{flushright}
\cleardoublepage

\thispagestyle{empty}
\vspace*{1.1in}
\begin{flushleft}
  {\color{accent}\rule{0.9in}{2.5pt}}\par\vspace{1.4ex}
  {\fontsize{32}{38}\selectfont\bfseries Machine\\[0.15ex] Learning\par}
  \vspace{2.2ex}
  {\Large\itshape From linear regression to convolutional\\neural networks, built from scratch in
    C++\par}
  \vspace{1.2in}
  {\large Erik Pihl\par}
  \vfill
  {\small\scshape\lsstyle qrtech academy\par}
  \vspace{0.6ex}
  {\small\color{muted} \bookedition\par}
\end{flushleft}
\newpage

\thispagestyle{empty}
\vspace*{\fill}
{\small
\noindent\textit{Machine Learning}\\
Copyright \textcopyright\ 2026 QRTECH Academy.

\medskip
\noindent The text and figures of this book, and of the course it is typeset from, are licensed
under the Creative Commons Attribution-NonCommercial-ShareAlike 4.0 International License
(CC BY-NC-SA 4.0). You may copy and redistribute them, and adapt them, for any non-commercial
purpose, provided you give appropriate credit and share what you make under the same license. To
view a copy of the license, visit \url{https://creativecommons.org/licenses/by-nc-sa/4.0/}.

\medskip
\noindent The code in this book may also be used under the MIT License, which is the license of all
the code in the companion repository. The repository holds the lecture material, the demo programs
and their reference implementations, the test suites, the finished network, and the solutions to
the exercises:

\medskip
\noindent\url{https://github.com/qrtech-academy/machine-learning}

\medskip
\noindent The C++ in this book is C++17 throughout, compiled with GCC, and uses nothing outside the
standard library: every model, layer and network in it is written from scratch. Arm Cortex-M0 and
Cortex-M0+ appear in the text as examples of processors without a hardware floating-point unit;
they are trademarks of their respective owners.

\medskip
\noindent Set in TeX Gyre Pagella and DejaVu Sans Mono with Lua\LaTeX.

\medskip
\noindent\textcolor{muted}{Typeset from the ten lectures of the course, their exercises, and the
two self-assessment papers. This edition: \bookedition.}
\par}
\cleardoublepage
